% Generated draft: requires mathematical and source review.
\documentclass[11pt]{article}
\usepackage[margin=25mm]{geometry}
\usepackage{fontspec}
\setmainfont{Latin Modern Roman}
\usepackage{amsmath,amssymb,mathtools,unicode-math}
\setmathfont{Latin Modern Math}
\usepackage{enumitem,needspace,xurl,hyperref}
\setlength{\parindent}{0pt}
\setlength{\parskip}{5pt}
\setlength{\emergencystretch}{3em}
\newcommand{\N}{\mathbb{N}}
\newcommand{\Z}{\mathbb{Z}}
\newcommand{\Q}{\mathbb{Q}}
\newcommand{\R}{\mathbb{R}}
\newcommand{\C}{\mathbb{C}}
\newcommand{\F}{\mathbb{F}}
\newcommand{\A}{\mathbb{A}}
\newcommand{\PP}{\mathbb{P}}
\newcommand{\E}{\mathbb{E}}
\newcommand{\eps}{\varepsilon}
\newcommand{\dd}{\,\mathrm{d}}
\begin{document}
\section*{UnsolvedMath ID: 2802}
\textbf{Source problem code:} catalogue entry
\subsection{Definitions and mathematical statement}
For odd integers $r\geq3$, put $q_r=e^{2\pi i/r}$. Then $q_r^2$ is also a primitive root of odd order $r$. Let $\mathrm{TV}(M;q)$ and $\mathrm{WRT}(M;q)$ denote the Turaev--Viro and Witten--Reshetikhin--Turaev invariants with the source's root-of-unity conventions, referenced to [TV92], [RT91], and [CY18]. For a hyperbolic $3$-manifold $M$, let $\operatorname{Vol}(M)$ be its hyperbolic volume and $\operatorname{CS}(M)$ its Chern--Simons invariant, with normalization and lift still to be specified for a complex equality.

The two Chen--Yang requests in the source are:
\begin{itemize}
\item (a) For every hyperbolic $3$-manifold $M$ in the intended invariant domain,
\[\lim_{\substack{r\to\infty\\r\ \mathrm{odd}}}\frac1r\log\bigl(\mathrm{TV}(M;q_r)\bigr)=\frac{\operatorname{Vol}(M)}{2\pi}.\]
The source explicitly allows closed manifolds, cusped manifolds, and manifolds with totally geodesic boundary. Its exact finite-volume/compact-core and boundary-state conventions are not given here; these must be fixed rather than silently restricting the phrase ``any hyperbolic $3$-manifold''.
\item (b) For every closed oriented hyperbolic $3$-manifold $M$,
\[\lim_{\substack{r\to\infty\\r\ \mathrm{odd}}}\frac1r\log\bigl(\mathrm{WRT}(M;q_r)\bigr)=\frac{\operatorname{Vol}(M)-i\operatorname{CS}(M)}{4\pi}.\]
The sign and coefficients are retained exactly, as is the full limit through odd levels rather than a chosen subsequence.
\end{itemize}
These formulas are targets, not already well-defined convention-free equalities. Part (b) needs a WRT normalization and phase, coherent logarithm choices along the levels, and a compatible representative of the Chern--Simons class. The principal logarithm alone has bounded imaginary part, whose division by $r$ tends to zero, so it cannot silently substitute for the branch information in the proposed complex-volume limit. The domain of the logarithms and handling of any zero invariant values also need to be part of a precise theorem. Part (a) likewise retains the source's logarithm rather than inserting an absolute value or changing to a limsup; positivity/nonvanishing conventions must follow its invariant definition.

The background credits both formulations to Qingtao Chen and Tian Yang [CY18]. The supplied triage checked 2026-08-17 reports (a) proved for many families, including broad cusped families, and records new gluing and Seifert-cobordism stability results of Detcherry--Kalfagianni--Marasinghe. It cites Wong--Yang for related relative-WRT asymptotics, not a universal proof of (b). It labels the programme partially solved / PARTIAL-PROGRESS and reports no general proof for every manifold in either requested class; the registry is open. These results and convention caveats are preserved without replacing either universal target by a verified-family or relative variant.
\subsection{Short English statement}
Do odd-level Turaev--Viro invariants recover hyperbolic volume, and do closed-manifold WRT invariants recover complex volume with the stated Chern--Simons phase? The source reports many established families but no universal proof, and the complex logarithm, phase, and normalization conventions remain essential to the formulation.
\subsection{Sources}
\begin{itemize}
\item Local UnsolvedMath registry, permanent ID 2802, source code catalogue entry, ``Kirby Problem 3.4''; supplied Chen--Yang formulas and literature notes.
\item [TV92] V. G. Turaev and O. Ya. Viro, ``State sum invariants of 3-manifolds and quantum 6j-symbols,'' \emph{Topology} 31(4) (1992), 865--902, doi:10.1016/0040-9383(92)90015-A.
\item [RT91] N. Reshetikhin and V. G. Turaev, ``Invariants of 3-manifolds via link polynomials and quantum groups,'' \emph{Inventiones Mathematicae} 103(3) (1991), 547--597, doi:10.1007/BF01239527.
\item [CY18] Qingtao Chen and Tian Yang, ``Volume conjectures for the Reshetikhin-Turaev and the Turaev-Viro invariants,'' \emph{Quantum Topology} 9(3) (2018), 419--460: \url{https://doi.org/10.4171/QT/111}. Supplied original formulation and evidence for both conjectures.
\item [DKM25] Renaud Detcherry, Efstratia Kalfagianni, and Shashini Marasinghe, ``Seifert cobordisms and the Chen-Yang volume conjecture'' (2025), arXiv:2505.01546: \url{https://arxiv.org/abs/2505.01546}. Supplied gluing and closure results for families, not universal resolution.
\item [WY23] Ka Ho Wong and Tian Yang, ``Relative Reshetikhin-Turaev invariants, hyperbolic cone metrics and discrete Fourier transforms I,'' \emph{Communications in Mathematical Physics} 400(2) (2023), 1019--1070, doi:10.1007/s00220-022-04613-5.
\item R. I. Baykur, R. C. Kirby, and D. Ruberman (eds.), \emph{K3: A New Problem List in Low-Dimensional Topology}, AMS 295 (2026), Problem 3.4: \url{https://math.berkeley.edu/sites/default/files/surv-295-ruberman-watermarked-author-pdf.pdf}. Supplied authoritative synthesis of known cases and the relative-WRT direction.
\item Supplied literature triage checked 2026-08-17, PARTIAL-PROGRESS; it explicitly preserves phase/logarithm/Chern--Simons caveats and cautions that dated triage does not exclude later or unindexed results.
\end{itemize}
\end{document}
